\section[Cote 120 : écraser une partie modérée sur un corps]{Cote 120 : écraser une partie modérée sur un corps\protect\verifie{Folder120}}\label{sec:120}

La cote 120, \emph{Topologie modérée}, vingt-quatre pages que l'inventaire date \q{à partir de 1973 --- à partir de 1978}, cherche des axiomes pour les parties \q{modérées} des puissances d'un ensemble $R$. La page~2 pose les axiomes (M1) à (M4) sur des familles $\mathcal M_n$ de parties de $R^n$, que la page~4 transcrit pour des ensembles d'indices finis quelconques ; la page~5 appelle modérée une application dont le graphe est une partie modérée. L'axiome (M6) de la page~12 demande que toute partie modérée $Y$ d'un espace modéré $X$ soit une fibre $\bigcap_\alpha f_\alpha^{-1}(e_\alpha)$ d'un nombre fini de fonctions modérées, les $e_\alpha$ étant des points modérés, \q{de plus, on peut supposer} que les $f_\alpha$ séparent les points de $X \smallsetminus Y$ : c'est dire que le quotient $X/Y$, où $Y$ est écrasé en un point, existe au sens strict. La section~(VI) suppose $R$ un corps et ajoute l'axiome (M11) : $0$ et $1$ sont des points modérés, $x - y$ et $xy$ sont des fonctions modérées sur toute partie modérée de $R^2$, et les inverses d'une partie modérée de $R^*$ sont contenus dans une partie modérée (\lot{120}{1}{20}). Les pages~22 et~23 en tirent que la seconde moitié de (M6) découle de la première. La page~22 se ramène à $e_\alpha = 0$ par translation et introduit les produits $f_\alpha g_\beta$. La page~23 vérifie la séparation et énonce la proposition \q{Si on a M1 : M7 et M11 \ldots}, dont la \q{Dém :} reste vide (\lot{120}{2}{22}). La page écrit \q{$e_\alpha \in R$} et \q{sépare les pts de $X$} ; la lecture corrige en $e_\alpha \in R_0$ et en $X \smallsetminus Y$.

\noindent\emph{Conventions.} Pour un ensemble fini $I$, $\mathcal M_I$ est un ensemble de parties de $R^I$ ; pour $u : I \to J$, $u^* : R^J \to R^I$ est $x \mapsto x \circ u$. Les axiomes, sous la forme (M$'$1) à (M$'$4) de la page~4, sont : (M$'$1) $u^*(A) \in \mathcal M_I$ pour $A \in \mathcal M_J$ ; (M$'$2) $(u^*)^{-1}(B) \in \mathcal M_J$ pour $u$ surjective et $B \in \mathcal M_I$ ; (M$'$3) $(u^*)^{-1}(B) \cap A \in \mathcal M_J$ ; (M$'$4) $A \times B \in \mathcal M_{I \sqcup J}$ pour $A \in \mathcal M_I$, $B \in \mathcal M_J$. Le graphe d'une famille $F : R^I \to R^K$ au-dessus de $X \subset R^I$ est $\Gamma_F = \{(x, F(x)) \mid x \in X\} \subset R^{I \sqcup K}$. Une fonction $f : X \to R$ est modérée si son graphe est dans $\mathcal M_{I \sqcup \{\ast\}}$, et $R_0$ est l'ensemble des $c \in R$ tels que $\{c\} \in \mathcal M_{\{\ast\}}$. Un espace modéré est ici une partie modérée $X \in \mathcal M_I$, et des fonctions modérées $g_\beta$ le plongent si $x \mapsto (g_\beta(x))_\beta$ est injective sur $X$.

\begin{lemme}[pages 5 et 20]\label{lem:120-anneau}
Soient $X \in \mathcal M_I$ et $f$, $g$ des fonctions modérées sur $X$.
\begin{enumerate}
\item Le graphe de $(f, g) : X \to R^2$ est modéré, ainsi que son image.
\item Si $F : R^I \to R^K$ a un graphe modéré au-dessus de $X$ et si $h$ est une fonction modérée sur $F(X)$, alors $h \circ F$ est modérée sur $X$.
\item Si $R$ est un anneau commutatif vérifiant (M11)~b), $f - g$ et $fg$ sont modérées. Les constantes $c \in R_0$ et les coordonnées $x \mapsto x_i$ sont modérées sur $X$.
\end{enumerate}
Seuls servent (M$'$1), (M$'$2), (M$'$4) et (M11)~b).
\end{lemme}

\begin{proof}
1. Soit $w : (I \sqcup \{\ast\}) \sqcup (I \sqcup \{\ast\}) \to I \sqcup \{0, 1\}$ l'application qui envoie les deux copies de $I$ sur $I$, la première étoile sur $0$ et la seconde sur $1$ ; elle est surjective. Un $y \in R^{I \sqcup \{0,1\}}$ a son image $w^*(y)$ dans $\Gamma_f \times \Gamma_g$ si et seulement si $y|_I \in X$, $y_0 = f(y|_I)$ et $y_1 = g(y|_I)$ : la préimage de $\Gamma_f \times \Gamma_g$, modérée par (M$'$4) puis (M$'$2), est le graphe de $(f, g)$. Plus généralement, l'image $F(X)$ d'une famille de graphe modéré est la projection de ce graphe sur $R^K$, modérée par (M$'$1).

2. On forme $\Gamma_F \times \Gamma_h \subset R^{(I \sqcup K) \sqcup (K \sqcup \{\ast\})}$, modéré par (M$'$4), et on le tire en arrière par la surjection $(I \sqcup K) \sqcup (K \sqcup \{\ast\}) \to (I \sqcup K) \sqcup \{\ast\}$ qui identifie les deux copies de $K$ : on obtient, par (M$'$2), l'ensemble des $(x, F(x), h(F(x)))$, $x \in X$. Sa projection sur $R^{I \sqcup \{\ast\}}$, modérée par (M$'$1), est le graphe de $h \circ F$.

3. On applique le point~2 à $F = (f, g)$, dont l'image est modérée par le point~1, et à $h$ égale à $(z_0, z_1) \mapsto z_0 - z_1$ ou $z_0 z_1$, modérée sur cette image par (M11)~b). Le graphe de la constante $c$ est $X \times \{c\}$, modéré par (M$'$4) ; celui de la coordonnée $x_i$ est l'image de $X$ par $u^*$, où $u : I \sqcup \{\ast\} \to I$ est l'identité sur $I$ et envoie $\ast$ sur $i$, modérée par (M$'$1).
\end{proof}

\begin{lemme}[pages 22 et 23, l'argument]\label{lem:120-separation}
Soient $R$ un anneau sans diviseur de zéro, $X$ un ensemble, $(f_\alpha)_{\alpha \in A}$ et $(g_\beta)_{\beta \in B}$ des familles quelconques de fonctions $X \to R$. La famille formée des $f_\alpha$ et des $f_\alpha g_\beta$ a pour lieu des zéros communs celui des $f_\alpha$. Si $x$ n'est pas un zéro commun des $f_\alpha$ et si un $g_\beta$ sépare $x$ et $y$, cette famille sépare $x$ et $y$.
\end{lemme}

\begin{proof}
Un zéro commun des $f_\alpha$ annule les $f_\alpha g_\beta$. Pour la séparation, si un $f_\alpha$ sépare $x$ et $y$, c'est fini. Sinon $f_\alpha(x) = f_\alpha(y)$ pour tout $\alpha$ ; on choisit $\alpha$ avec $\lambda = f_\alpha(x) \neq 0$ et $\beta$ avec $g_\beta(x) \neq g_\beta(y)$. Alors $f_\alpha g_\beta(x) - f_\alpha g_\beta(y) = \lambda\,(g_\beta(x) - g_\beta(y))$ n'est pas nul, puisque $R$ n'a pas de diviseur de zéro.
\end{proof}

\begin{theoreme}[Grothendieck, pages 22 et 23, écrasement sur un corps\piste{120-collapse-separation-over-a-field}]\label{thm:120-ecrasement}
Supposons (M$'$1) à (M$'$4) et (M11), $R$ étant un corps. Soient $X \in \mathcal M_I$, des fonctions modérées $f_\alpha$ sur $X$ en nombre fini, des $e_\alpha \in R_0$, $Y = \bigcap_\alpha f_\alpha^{-1}(e_\alpha)$, et des fonctions modérées $g_\beta$ qui plongent $X$. Les fonctions $f_\alpha - e_\alpha$ et $(f_\alpha - e_\alpha)\,g_\beta$ sont modérées, ont $Y$ pour lieu des zéros communs, et séparent les points de $X \smallsetminus Y$. En particulier $Y$ est le lieu des zéros communs d'une famille finie de fonctions modérées qui sépare les points de $X \smallsetminus Y$ : la seconde moitié de (M6) découle de la première, et $X/Y$ est un espace modéré au sens strict.
\end{theoreme}

\begin{proof}
Par le lemme~\ref{lem:120-anneau}, la constante $e_\alpha$ est modérée parce que $e_\alpha \in R_0$, donc $f_\alpha - e_\alpha$ aussi, et les produits $(f_\alpha - e_\alpha)\,g_\beta$ le sont. Les zéros communs des $f_\alpha - e_\alpha$ forment $Y$. Pour $x \neq y$ dans $X \smallsetminus Y$, un $f_\alpha - e_\alpha$ ne s'annule pas en $x$, et comme les $g_\beta$ plongent $X$, l'un d'eux sépare $x$ et $y$ : le lemme~\ref{lem:120-separation} conclut. Pour l'énoncé d'existence, les coordonnées, modérées par le lemme~\ref{lem:120-anneau}, plongent $X$, et la famille est indexée par $A \sqcup (A \times I)$, qui est fini.
\end{proof}

\begin{remarque}
L'énoncé tient tel que la lecture le donne, et plus généralement. Ni (M$'$3), ni (M5) à (M7), ni les parties a) et c) de (M11), ni la finitude des familles ne servent ; de (M11), seule sert la partie b), sur un anneau commutatif sans diviseur de zéro. L'hypothèse d'intégrité ne peut être omise : sur $\mathbf Z/4\mathbf Z$, avec $f = 2$ et $g = \mathrm{id}$, la famille $2,\ 2x$ ne sépare pas $0$ et $2$, bien que $f(0) \neq 0$ et $g(0) \neq g(2)$. Les deux corrections de la lecture sont nécessaires. Sans $e_\alpha \in R_0$, la constante $e_\alpha$ n'est pas modérée et la translation n'est pas permise. Et une famille qui s'annule sur $Y$ ne sépare aucun couple de points de $Y$ : \q{sépare les points de $X$} ne peut s'entendre que de $X \smallsetminus Y$.
\end{remarque}

Ne sont pas démontrés ici : la catégorie $\Sigma$ des espaces modérés et ses limites finies (pages~5 à~7), l'existence des sommes amalgamées le long des plongements modérés sous (M5$'$), (M6) et (M7) (pages~12 et~13\piste{120-axiomatic-pushouts-along-embeddings}), les ind-espaces et la topologie des pages~14 à~19, et le programme de trivialité de la page~24.
